\section{Implementation and Architecture}\label{sect:ImplementationArchitecture}

The implementation of TrainHub follows the distinction established during the design phase: members and professionals use different interfaces, but both work on the same information. A workout assigned by a trainer, a message sent in chat, or a confirmed appointment must be visible from the other side without creating separate copies of the data. For this reason, the application was developed as a single Progressive Web App connected to a shared backend.

This chapter focuses on the structure that supports the final interface. Rather than describing the screens again, it explains how navigation, local state, persistent data, offline behaviour, and device features cooperate during the main user flows.

\subsection{Technology Stack}\label{subsect:TechnologyStack}

TrainHub is a client-side web application built with React and Vite. React provides the component model used throughout the interface, while Vite manages the development environment and creates the optimized production bundle. Material UI supplies the component foundation and theme system, and React Router maps the two role-specific navigation trees described in Section \ref{subsect:NavigationMap}.

Supabase provides authentication, the PostgreSQL database, Realtime updates, and the server-side function used to deliver push notifications. TanStack Query sits between the interface and the data functions: it manages remote state, keeps query results synchronized, and persists the useful part of the cache in IndexedDB. The main technologies, the versions the delivered build was compiled against, and their role in the project are summarized in Table \ref{tab:technologyStack}.

\begin{table}[H]
    \centering
    \small
    \begin{tabularx}{\textwidth}{@{}p{0.22\textwidth} p{0.11\textwidth} >{\raggedright\arraybackslash}X@{}}
        \toprule
        \textbf{Technology} & \textbf{Version} & \textbf{Use in TrainHub} \\
        \midrule
        React & 19.2 & Component-based interface and the hooks that hold screen-local state \\
        Vite & 8.1 & Development server, production build, and route-level code splitting \\
        Material UI & 9.2 & Responsive components, shared visual rules, typography, and theme customization \\
        React Router & 8.3 & Public routes, protected member and professional areas, and route-level error handling \\
        TanStack Query & 5.101 & Loading and mutation state, cache invalidation, offline queues, and persisted server data \\
        Supabase JS & 2.110 & User sessions, PostgreSQL access, Realtime channels, and stored procedure calls \\
        vite-plugin-pwa & 1.3 & Manifest injection and the build step that lists the files to precache \\
        Workbox & 7.4 & Precaching and navigation routing inside the hand-written service worker \\
        \bottomrule
    \end{tabularx}
    \caption{Main technologies used in the final application.}
    \label{tab:technologyStack}
\end{table}

The application is written in JavaScript and JSX. Keeping one language across the interface and data-access layer made the project easier to inspect during development, while the database remains responsible for rules that cannot safely depend on the browser.

\subsection{Application Structure}\label{subsect:ApplicationStructure}

The source code is organized by responsibility and by feature. Shared building blocks remain separate from domain-specific screens, while calls to persistent data are collected in plain asynchronous functions. Table \ref{tab:sourceStructure} shows the main directories.

\begin{table}[H]
    \centering
    \small
    \begin{tabularx}{\textwidth}{@{}p{0.24\textwidth} >{\raggedright\arraybackslash}X@{}}
        \toprule
        \textbf{Directory} & \textbf{Responsibility} \\
        \midrule
        \path{src/routes} & Public routes and the protected member and professional route trees \\
        \path{src/layouts} & The signed-in application shell and the centred layout used by the authentication screens \\
        \path{src/features} & Screens and feature-specific components, one folder per domain: authentication, home, agenda, workout, nutrition, trainer, clients, calendar, chat, notifications, progress, rewards, check-in, and profile \\
        \path{src/components} & Reusable interface elements such as page headers, navigation controls, status feedback, and shared cards \\
        \path{src/data} & Database reads and writes, grouped by application domain and independent from React \\
        \path{src/lib} & Supabase configuration, query keys, cache persistence, date utilities, identifiers, and other shared services \\
        \path{src/theme} & Material UI theme and the visual tokens exported from the Figma design \\
        \bottomrule
    \end{tabularx}
    \caption{Responsibilities of the main source directories.}
    \label{tab:sourceStructure}
\end{table}

The dependency direction is deliberately one-way. Modules in \path{src/data} know about Supabase and nothing about React; modules in \path{src/features} know about React and call \path{src/data}; components know about neither and render what they are given. A pure module such as the week calculation in \path{src/lib} can therefore be used by a data function and by a screen at the same time, which would not be possible if the shared code lived beside the screens that first needed it.

This organization also avoids placing database details inside every screen. A page requests the information it needs through a query, and the corresponding function in \path{src/data} translates that request into a Supabase operation. The same functions can therefore be used from different screens without duplicating filters, date conversions, or error handling. Authentication and Realtime subscriptions are the main exceptions: they are tied to the browser session and the component life cycle, so they remain close to the providers and hooks that own them.

\subsubsection{Application Shell and Routing}\label{subsubsect:ApplicationShellRouting}

At start-up, the root component installs the Material UI theme, the persisted query client, the authentication provider, and the router. The authentication provider first restores the Supabase session and then loads the associated profile. The profile identifies the user's role and assigned professional, information that is required by most protected flows.

Member routes use the \path{/m} prefix and professional routes use \path{/p}. The guard is the shell itself: each protected tree declares the role it expects, so a signed-in user who opens the other area is redirected to their own home page, while a visitor who is not signed in is sent to the login screen with the requested address remembered for after. Both trees are rendered inside that same shell, which contains the top bar, bottom navigation, offline feedback, membership banner, notification status, and the mini-player shown during an open workout. Screen modules are loaded only when their route is visited, so each section ships as its own chunk and the first view downloads a fraction of the application.

Detail pages return to their parent section through one shared header rather than through a back arrow drawn per screen, so a link opened from a notification still lands on a complete page with a valid way out. Unexpected route errors are handled by a shared error screen instead of leaving the application blank.

\subsubsection{Remote State and Mutations}\label{subsubsect:RemoteStateMutations}

Information read from Supabase is treated as remote state. Every query uses a shared key that identifies both the domain and the relevant user or record. When an operation changes data, related keys are invalidated together so that all affected views are refreshed. For example, completing a workout updates the current run, the weekly workout summary, the rewards area, and the professional's progress view without requiring those screens to know about each other.

Two defaults shape the whole offline behaviour, and they are set once for every query in the application.

\begin{lstlisting}[language=JavaScript,caption={Query defaults, in \texttt{src/lib/queryClient.js}.},label={lst:queryDefaults}]
queries: {
  gcTime: ONE_WEEK,
  staleTime: 1000 * 30,
  // Retrying while offline pauses the query instead of failing it, so
  // `isPending` would stay true forever and no screen could ever show
  // its offline state.
  retry: (failureCount) => navigator.onLine && failureCount < 2,
  refetchOnWindowFocus: false,
  // Serve cached data immediately and only then reach for the network.
  networkMode: 'offlineFirst',
}
\end{lstlisting}

The consequence is that a failed refresh is not the same as an absent answer. Screens decide their error state on whether data exists at all, never on the failure of the last request: an error page drawn over perfectly usable cached data would tell an offline member that the application is broken while it is doing exactly what it was built to do.

Writes follow the same pattern, but they need one thing more. The function a write executes is registered centrally, before the persisted cache is restored, rather than being passed at the call site. This matters for an offline mutation: after the application has been closed and reopened, the stored operation is found again by its key, and a restored write with no registered function is discarded without a message. Operations whose order matters share a scope, because paused writes are otherwise replayed in parallel and a set could reach the database before the run it belongs to.

\begin{lstlisting}[language=JavaScript,caption={A registered write, with the scope that keeps the workout writes in order.},label={lst:mutationDefault}]
queryClient.setMutationDefaults(mutationKeys.resumeRun, {
  mutationFn: resumeRun,
  onSuccess: (data, variables) => { /* update the open run in cache */ },
  scope: runScope,
  onSettled: () => {
    queryClient.invalidateQueries({ queryKey: queryPrefixes.runs })
  },
})
\end{lstlisting}

Some immediate actions are deliberately excluded from this queue. A badge token that expires after a minute, the redemption of a gym entry, and the registration of a push subscription cannot represent a real event if they are replayed much later. When the device is offline, these actions fail immediately and the interface explains that a connection is required.

\subsection{Backend and Data Model}\label{subsect:BackendDataModel}

The backend uses a relational PostgreSQL schema of twenty-one tables, because the principal entities are strongly connected. A member belongs to a professional, a plan contains sessions, a session contains exercises, and every completed set belongs to one specific workout run. Preserving these relations in the database makes it possible to enforce ownership and consistency independently of the interface.

The data model can be divided into the areas shown in Table \ref{tab:dataAreas}.

\begin{table}[H]
    \centering
    \small
    \begin{tabularx}{\textwidth}{@{}p{0.26\textwidth} >{\raggedright\arraybackslash}X@{}}
        \toprule
        \textbf{Area} & \textbf{Stored information} \\
        \midrule
        Accounts & Profiles, roles, professional assignment, specialty, and membership status \\
        Workouts & Exercise catalogue, versioned plans, sessions, prescribed exercises, workout runs, individual set logs, and rewards \\
        Nutrition & Versioned plans, reusable day types assigned to weekdays, and meals with their macronutrients \\
        Organization & Professional availability, appointment requests, confirmations, and cancellations \\
        Communication & One member--professional thread, messages, read receipts, and persistent notifications \\
        Facility access & Short-lived badge tokens and recorded check-ins \\
        Devices and configuration & Registered push endpoints and the server-side settings used to deliver them \\
        \bottomrule
    \end{tabularx}
    \caption{Main areas of the TrainHub data model.}
    \label{tab:dataAreas}
\end{table}

The schema is relational with one deliberate exception: the foods inside a meal are stored as a JSON document rather than as a further table. They are never queried on their own, never referenced from elsewhere, and are always read and written together with the meal that contains them, so a fourth level of joins would have bought nothing.

Workout and nutrition plans are not overwritten when they are edited. The replacement records a link to the plan it supersedes, and a partial unique index allows any plan to be superseded only once, while workout runs continue to reference the plan and session that were actually used. This prevents a later change to the programme from rewriting the meaning of an older result. The current plan is simply the most recent one created for that member.

A similar distinction is used between a workout session and a workout run. The session describes what should be performed; the run records one attempt, including start time, pauses, logged sets, final percentage, outcome, and the member's closing note. Weekly progress is derived from these runs rather than stored as a permanent status on the session, so the same session can be trained more than once and a partial attempt remains different from a completed one. The session table still carries the status column of the earlier design; nothing reads it and nothing decides on it, and it was left in place because dropping a column is a destructive migration that would have bought only tidiness.

\subsubsection{Authorization and Consistent Writes}\label{subsubsect:AuthorizationConsistentWrites}

The browser connects to Supabase with the publishable project key. That key is compiled into the JavaScript bundle and is public by construction: it identifies the project and grants nothing on its own. Every request carries the authenticated user identity, and PostgreSQL answers it in two stages that are easy to confuse. It first checks the table privileges, and only then evaluates the Row Level Security policies. A table with impeccable policies and no privilege answers \texttt{permission denied} to everybody, while a policy whose condition never mentions the current user identity is readable by anyone holding the key. Both failures are silent from the application's point of view, and an early version of this project shipped the second one.

The rule that separates a member from a professional is written once and reused by every policy that needs it.

\begin{lstlisting}[language=SQL,caption={The predicate that encodes ownership, and one of the policies built on it.},label={lst:ownsMember}]
create function owns_member(target uuid) returns boolean
language sql security definer stable set search_path = '' as $$
  select target = auth.uid()
      or exists (
           select 1 from public.profiles
           where id = target and assigned_pro_id = auth.uid()
         );
$$;

create policy workout_runs_select on workout_runs
  for select using (owns_member(member_id));
\end{lstlisting}

The empty search path is a security measure rather than a formality: inside a function that runs with the definer's privileges, an unqualified table name would be resolved against the caller's temporary schema first, and a signed-in user controls that schema. One table, holding the settings used to deliver notifications, has neither a policy nor a privilege and is reachable only from inside the database.

Simple operations are protected by these policies alone. Everything that changes data goes through stored procedures instead, and the direct write privileges on those tables were revoked. Creating a workout plan inserts the plan, every drafted session, and all their exercises as a single transaction. Closing a workout computes the result from the sets that were actually logged and writes the reward in the same operation. If any part fails, the database is not left holding half a plan or a summary that describes work nobody did.

Client-generated identifiers make an exact retry safe after a lost connection: the second attempt lands on the same row instead of creating a second one. The same procedures compare the content behind a reused identifier and reject a collision rather than accepting it silently. Plan creation carries one further check: the caller sends the identifier of the version it believes it is replacing, and the database refuses the write if a newer plan has appeared in the meantime, so a member and their professional editing at the same time cannot overwrite each other unnoticed.

Membership rules are enforced at this level as well as in the interface. An inactive member cannot create a plan, request an appointment, enter through the QR badge, or earn a workout reward, although the run itself is still recorded. A professional may change the status only for a client assigned to them. Keeping these checks in the backend prevents a modified browser request from bypassing the disabled buttons shown on screen.

\subsection{Progressive Web App and Offline Behaviour}\label{subsect:PWAOffline}

TrainHub includes a web app manifest with its name, description, icons, theme and background colours, portrait orientation, and standalone display mode. The two colours are the values the Figma tokens define, so the installed icon and splash screen cannot drift away from the interface they introduce.

The service worker is written by hand rather than generated. The plugin runs in the mode that only injects the list of files to precache, because the same file also has to carry the push and notification-click handlers, and a generated worker has nowhere clean to put them.

\begin{lstlisting}[language=JavaScript,caption={The service-worker build configuration, in \texttt{vite.config.js}.},label={lst:injectManifest}]
VitePWA({
  registerType: 'autoUpdate',
  strategies: 'injectManifest',
  srcDir: 'src',
  filename: 'sw.js',
  injectManifest: {
    // Inter ships seven unicode-range subsets. Precaching all of them
    // would fetch glyphs this application never renders.
    globPatterns: ['**/*.{js,css,html}', 'assets/inter-latin*.woff2'],
  },
})
\end{lstlisting}

The precache holds the entry page, every route chunk, the stylesheet, the icon set, and only the Latin subsets of the typeface. Navigation requests fall back to the stored application shell, so an installed copy starts even when the network is unavailable.

Application data is handled separately, and the service worker caches none of it. A shared cache holding one user's rows would survive a sign-out on a shared device, which is a leak rather than an optimisation. Successful query results are instead persisted in IndexedDB for at most a week, keyed to the build that wrote them. A previously loaded workout plan or appointment list therefore remains visible offline, while the interface reports that it may not be current. Signing out clears the local query cache and removes the push subscription registered for that device.

Mutations that remain meaningful after reconnection are stored alongside the cache. Each queued operation keeps the client-generated identifier and the timestamp of the original action, so a replay does not turn an old event into a new one. When connectivity returns, the pending operations resume in scope order and the affected queries refresh. A migration step runs before anything is replayed and discards operations written by an incompatible earlier build, telling the user once that this happened rather than sending an obsolete request to a database that has moved on.

Offline support is therefore selective. Browsing recently loaded information and recording recoverable changes continue to work; actions that need a current server decision, such as authentication, badge validation, or enabling push notifications, require a connection. This boundary keeps the offline experience useful without pretending that every function can work in isolation.

\subsubsection{Installability and Delivery}\label{subsubsect:InstallabilityDelivery}

A browser offers installation only over a secure origin, so the application is served through an HTTPS tunnel during development and demonstration. The tunnel uses a fixed domain rather than the address a free account generates on each restart, and the reason is not convenience. A Web Push subscription is bound to the origin that created it, and the authentication provider matches its redirect addresses against a list. A changing origin would invalidate every registered device and break the password-recovery link every time the tunnel came up.

The service worker does not exist under the development server, which is deliberate: an application that reloads on every keystroke and a cache that answers before the network are difficult to reconcile. Installability, offline start-up, and push delivery are verified against a production build served locally, which is also the configuration used on the day of the demonstration.

\subsection{Dates, Weeks, and Time Zones}\label{subsect:DatesWeeks}

Almost every screen answers a question about a day: which sessions belong to this week, which meals belong to this weekday, whether a membership has expired, whether a badge has already been redeemed today. Dates were one of the few areas where a small mistake produced a visible defect that was hard to reproduce, and the rules that emerged are worth stating.

A date written as a plain calendar string is parsed as midnight UTC, which renders as the previous day for every user west of Greenwich. All conversions pass through two helpers that compute in the local frame instead. Days are never subtracted as milliseconds, because one day a year lasts twenty-three hours and another twenty-five. Week membership is decided on the local day of a run rather than on the raw timestamp: a session started at half past eleven on a Sunday evening carries a UTC timestamp already dated Monday, and comparing the strings would file it under a week that has not started.

\begin{lstlisting}[language=JavaScript,caption={Deriving a session's weekly state from its runs, in \texttt{src/lib/week.js}.},label={lst:runStatus}]
export function runStatusOf(runs, weekStartISO) {
  const open = (runs ?? []).find((run) => !run.ended_at)
  if (open) return { status: 'in_progress', run: open }

  const thisWeek = (runs ?? [])
    .filter((r) => r.outcome === 'completed' || r.outcome === 'partial')
    .filter((r) => localDayISO(r.started_at) >= weekStartISO)
    .sort((a, b) => Date.parse(b.started_at) - Date.parse(a.started_at))

  if (thisWeek.length === 0) return { status: 'todo', run: null }
  return { status: thisWeek[0].outcome, run: thisWeek[0] }
}
\end{lstlisting}

The same question arises on the server, where the database clock is not the member's. The reward for a gym check-in is stamped with the Italian calendar day rather than with the server's day, so a scan a few minutes after midnight cannot pay twice, nor be discarded as a duplicate of the evening before.

\subsection{A Workout from Start to Finish}\label{subsect:WorkoutEndToEnd}

The workout flow shows how the layers cooperate during an everyday task:

\begin{enumerate}
    \item The workout page requests the member's current plan and the runs recorded during the current week. TanStack Query first makes any valid persisted result available, then checks Supabase when the network permits it.
    \item When the member starts a session, the application generates the run identifier and asks a protected database function to open it. The database verifies the member, and a partial unique index prevents a second open run from existing.
    \item Every completed set is associated with both the run and the prescribed exercise. The interface updates immediately, while the shared mutation scope preserves the sequence if the connection is interrupted.
    \item Elapsed time is derived from stored timestamps and accumulated pause time. It does not depend on a counter that would be lost when the screen closes, so the mini-player can follow the run across the member area.
    \item At the end, the database compares the valid logged sets with the prescribed targets, records a completed, partial, or abandoned outcome, calculates the percentage, and awards points when the membership permits it.
    \item The returned result invalidates the relevant queries. The member sees the run summary and the updated weekly progress, while the professional can open the same run from the client's progress area and read its note.
\end{enumerate}

Step five is where the application refuses to trust the client. The percentage and the points are recomputed from the rows that exist, not from what the screen believed it had recorded.

\begin{lstlisting}[language=SQL,caption={Closing a run: the result is derived from the set logs.},label={lst:closeRun}]
select coalesce(sum(least(coalesce(logged.amount, 0), item.target_sets)), 0)
into v_done
from public.session_exercises item
left join (
  select session_exercise_id, count(*)::int amount from public.set_logs
  where run_id = p_run_id group by session_exercise_id
) logged on logged.session_exercise_id = item.id
where item.session_id = v_row.session_id;

v_pct := case when v_target = 0 then 0
              else floor(100.0 * v_done / v_target) end;
v_points := case when v_outcome = 'abandoned' or v_target = 0 then 0
                 else floor(30.0 * v_done / v_target) end;
\end{lstlisting}

The visible percentage consequently represents performed work rather than a visit to the session page. If a session contains four exercises and the member records only one, the weekly summary remains partial. This distinction is essential for both the member's feedback and the professional's interpretation of progress.

\subsection{Realtime, Push, and Device Features}\label{subsect:RealtimePushDevice}

Chat uses Supabase Realtime to receive inserted messages and updated read receipts. The query cache remains the source of truth: Realtime events add or replace matching rows, and an identifier already present is ignored, which is what stops a message from appearing twice while the sender's own copy is still on screen. A periodic refresh provides a recovery path when a mobile network silently interrupts the socket, so the conversation stays usable even if the live connection is temporarily unavailable.

Notifications are stored in the database before they are delivered to a device. The in-app notification centre keeps its history even when Web Push is unsupported or permission is denied. Registered browsers receive the same event through an Edge Function that signs the message with VAPID keys; that function is called by the database itself when the row is written, not by the application. The service worker displays the notification outside the application and, when it is tapped, reuses an already open TrainHub window instead of opening a second one. Endpoints that a push service reports as gone are deleted, so a browser that discarded its subscription is not contacted again.

Facility access uses two native browser capabilities. The member's page draws a QR code from a short-lived random token, while the professional's scanner requests camera access and decodes the image locally, on a canvas, without a frame leaving the device. Redemption is performed by the database, which verifies that the caller is a professional, that the token exists, that it has not already been used, that it has not expired, and that the member's subscription allows entry, before recording the check-in and its reward. The token is drawn from an alphabet without visually ambiguous characters, so it can be read aloud or typed by hand when a camera fails.

These functions enhance the application without becoming dependencies for the remaining flows. A user who refuses notification or camera permission can still use workouts, nutrition, appointments, and chat; only the feature that requires that permission is unavailable.

\subsection{Working Within the Platform}\label{subsect:PlatformLimits}

A Progressive Web App runs inside the browser's rules, and several of those rules only became visible on a real device. They are recorded here because they shaped the implementation more than any library choice did.

The clearest case was the rest timer. Its chime was built on the Web Audio API and was silent on iPhone for two independent reasons. The audio context was created inside the timer callback, ninety seconds after the tap that started the rest, and a context created outside a user gesture is born suspended and never sounds. Correcting that was not enough: Safari routes Web Audio through the ambient category, which the hardware silent switch mutes, and most phones in a gym have that switch on. The timer now uses an ordinary audio element, unlocked inside the tap that starts the rest, because that path ignores the switch. Making it audible then created a third problem, since playing it interrupted the member's own music; the unlock is now performed muted, and the sound is declared transient, which is the category the system's own timer uses and which lowers other audio instead of stopping it.

Three further absences had to be designed around rather than solved. The Vibration API does not exist in Safari on iOS, so touch feedback cannot carry information. The web has no scheduled local notification at all, so nothing can be arranged to fire while the application is closed. Background timers are throttled by every browser and frozen by iOS, which is the reason elapsed time is recomputed from stored timestamps rather than counted. For the same reason the end of a rest is also announced in text through a live region, so a member who never hears the chime is never left waiting for it.

\subsection{Validation and Verification}\label{subsect:ValidationVerification}

Validation takes place in both the interface and the backend. Forms reject incomplete or inconsistent values before submission, while database constraints and protected functions repeat the checks that concern ownership or stored data. Screens distinguish loading, empty, failed, and offline states; a failed background refresh does not discard a valid cached result already shown to the user.

No test runner was adopted. Pure logic that is worth protecting instead ships with an assertion script beside it, run directly with Node: seventeen of them cover dates and week boundaries, the workout timer, run summaries and reward points, the workout and nutrition drafts, QR scan outcomes, membership state, push key encoding, the calendar grid, and the cache migration. They are cheap enough to run after every change and they fail loudly, which is what a small team actually needs from a test.

Both \texttt{npm run lint} and \texttt{npm run build} are required before a change is considered finished, and the second is not a formality: ESLint does not resolve module paths, so an icon that the installed library does not export passes linting and breaks the build.

The database is verified from three directions, because no single one is sufficient. A verification script checks the schema, the privileges, the policies, the protected functions, and a set of integrity counts that must all read zero, such as a run whose percentage disagrees with its own set logs. It runs with administrative rights, however, and consequently bypasses Row Level Security entirely: it can prove that a policy exists, not that the policy is correct. A second script asks the database the same question an attacker would, holding nothing but the publishable key from the bundle, and requires every table to answer with no rows. A third signs in as each demo role and confirms that the tables the application needs are genuinely reachable, which catches the opposite failure of a policy so strict that a legitimate screen renders empty.

Manual testing then followed complete scenarios with both roles, from plan creation and workout progress to appointments, chat, notifications, access badges, and membership status. Installation, offline start-up, camera access, and push delivery were checked from a production preview on real devices, where these browser capabilities can be tested properly.
